% =====================================================================
%  COURS DE PHYSIQUE — FICHE 6
%  Ondes (propagation, onde sinusoïdale, ondes stationnaires,
%         diffraction & interférences)
%  Figures TikZ tracées « from scratch ». Sans section exercices.
%  Compilation : pdflatex (2 passes pour la table des matières).
% =====================================================================
\documentclass[11pt,a4paper]{article}

% ---------- Encodage & langue ----------
\usepackage[utf8]{inputenc}
\usepackage[T1]{fontenc}
\usepackage{lmodern}
\usepackage[french]{babel}

% ---------- Mathématiques ----------
\usepackage{amsmath,amssymb,mathtools}
\usepackage{icomma}

% ---------- Unités ----------
\usepackage{siunitx}
\sisetup{output-decimal-marker={,},per-mode=symbol}

% ---------- Couleurs, boîtes, graphismes ----------
\usepackage{xcolor}
\usepackage{colortbl}
\usepackage{tikz}
\usepackage{pgfplots}
\usepackage[most]{tcolorbox}
\usepackage{enumitem}
\usepackage{array}
\usepackage{booktabs}
\usepackage{multicol}
\usepackage{caption}
\usepackage{fancyhdr}
\usepackage[hidelinks]{hyperref}

\usetikzlibrary{arrows.meta,positioning,calc,decorations.pathmorphing,
  decorations.markings,angles,quotes,patterns,backgrounds,
  shapes.geometric,shapes.misc,fadings,intersections,fit}
\pgfplotsset{compat=1.18}

% ---------- Géométrie de page ----------
\usepackage[a4paper,margin=2.2cm,headheight=26pt,headsep=10pt]{geometry}

% =====================================================================
%  PALETTE DE COULEURS
% =====================================================================
\definecolor{primary}{RGB}{0,151,169}
\definecolor{primarydark}{RGB}{0,88,101}
\definecolor{defcol}{RGB}{0,114,178}
\definecolor{thmcol}{RGB}{0,140,120}
\definecolor{formulacol}{RGB}{198,93,9}
\definecolor{remarkcol}{RGB}{120,94,200}
\definecolor{attncol}{RGB}{200,40,55}
\definecolor{excol}{RGB}{0,135,90}
\definecolor{methcol}{RGB}{90,90,100}
\definecolor{ondeA}{RGB}{0,114,178}   % bleu  (onde 1 / onde incidente)
\definecolor{ondeB}{RGB}{220,110,30}  % orange (onde 2 / onde réfléchie)

% =====================================================================
%  STYLE COMMUN DES BOÎTES
% =====================================================================
\tcbset{
  boxbase/.style={enhanced,breakable,boxrule=0.9pt,arc=2.5pt,
     left=3mm,right=3mm,top=2.3mm,bottom=2.3mm,
     fonttitle=\bfseries,coltitle=white,
     toptitle=0.6mm,bottomtitle=0.6mm},
}
\newtcolorbox[auto counter,number within=section]{definition}[1][]{%
  boxbase,colback=defcol!5,colframe=defcol,
  title={\textsc{Définition}~\thetcbcounter\ \ifx\relax#1\relax\relax\else\textmd{--- \itshape #1}\fi}}
\newtcolorbox[auto counter,number within=section]{theoreme}[1][]{%
  boxbase,colback=thmcol!6,colframe=thmcol,
  title={\textsc{Loi / Propriété}~\thetcbcounter\ \ifx\relax#1\relax\relax\else\textmd{--- \itshape #1}\fi}}
\newtcolorbox[auto counter,number within=section]{formule}[1][]{%
  boxbase,colback=formulacol!6,colframe=formulacol,
  title={\textsc{Formule}~\thetcbcounter\ \ifx\relax#1\relax\relax\else\textmd{--- \itshape #1}\fi}}
\newtcolorbox{remarque}[1][]{%
  boxbase,colback=remarkcol!6,colframe=remarkcol,
  title={\textsc{Remarque}\ifx\relax#1\relax\relax\else\textmd{ : \itshape #1}\fi}}
\newtcolorbox{attention}[1][]{%
  boxbase,colback=attncol!6,colframe=attncol,
  title={\textsc{Attention}\ifx\relax#1\relax\relax\else\textmd{ : \itshape #1}\fi}}
\newtcolorbox{exemple}[1][]{%
  boxbase,colback=excol!5,colframe=excol,
  title={\textsc{Exemple}\ifx\relax#1\relax\relax\else\textmd{ : \itshape #1}\fi}}
\newtcolorbox{methode}[1][]{%
  boxbase,colback=methcol!7,colframe=methcol,sharp corners,
  title={\textsc{Méthode}\ifx\relax#1\relax\relax\else\textmd{ : \itshape #1}\fi}}

% =====================================================================
%  ENVIRONNEMENT « où : »
% =====================================================================
\newenvironment{ou}{%
  \par\smallskip\noindent\textbf{où :}\nopagebreak
  \begin{itemize}[leftmargin=1.8em,topsep=2pt,itemsep=1.5pt,parsep=0pt]}%
  {\end{itemize}}

% =====================================================================
%  EN-TÊTES ET PIEDS
% =====================================================================
\pagestyle{fancy}
\fancyhf{}
\fancyhead[L]{\small\textcolor{primarydark}{\textbf{Fiche 6} — Ondes}}
\fancyhead[R]{\small\textcolor{primarydark}{\textsc{Physique}}}
\fancyfoot[C]{\small\thepage}
\renewcommand{\headrulewidth}{0.8pt}
\renewcommand{\headrule}{\hbox to\headwidth{\color{primary}\leaders\hrule height \headrulewidth\hfill}}

\usepackage{titlesec}
\titleformat{\section}{\Large\bfseries\color{primarydark}}{\thesection}{0.6em}{}
\titleformat{\subsection}{\large\bfseries\color{primary}}{\thesubsection}{0.6em}{}
\titleformat{\subsubsection}{\normalsize\bfseries\color{primary!80!black}}{\thesubsubsection}{0.6em}{}

% Raccourcis maths
\newcommand{\vv}[1]{\vec{#1}}
\newcommand{\dd}{\mathrm{d}}

% =====================================================================
\begin{document}
\sloppy

% ---------------------------------------------------------------------
%  PAGE DE TITRE
% ---------------------------------------------------------------------
\begingroup
\thispagestyle{empty}
\centering
\vspace*{1.4cm}

\begin{tikzpicture}
\node[rectangle,rounded corners=8pt,fill=primary,text=white,
      minimum width=15.5cm,minimum height=2.6cm,align=center,inner sep=10pt]
  {\Huge\bfseries Ondes\\[6pt]
   \large\bfseries Propagation, ondes sinusoïdales, ondes stationnaires,\\[2pt]
   diffraction \& interférences};
\end{tikzpicture}

\vspace{0.4cm}
{\large\color{primarydark}\textsc{Cours de physique — Fiche n\textsuperscript{o}\,6}}

\vspace{0.9cm}

% --- Illustration de couverture : une onde sinusoïdale + longueur d'onde ---
\begin{tikzpicture}[scale=1.0]
  \def\A{1.0}\def\L{2.4}
  % axe horizontal
  \draw[->,>=Stealth,black!70] (-0.4,0) -- (10.4,0) node[right]{$x$};
  \draw[->,>=Stealth,black!70] (0,-1.5) -- (0,1.7) node[above]{$y$};
  % sinusoïde (deux longueurs d'onde et demie)
  \draw[primary,line width=1.4pt,smooth,samples=220,domain=0:9.6]
        plot (\x,{\A*sin( \x*(360/\L) )});
  % marquage de la longueur d'onde entre deux crêtes
  \coordinate (c1) at ({0.25*\L},\A);
  \coordinate (c2) at ({1.25*\L},\A);
  \draw[<->,>=Stealth,primarydark,thick] (c1)++(0,0.35) -- ($(c2)+(0,0.35)$)
        node[midway,above,fill=white,inner sep=1pt]{$\lambda$};
  \draw[primarydark,dashed,line width=0.5pt] (c1) -- ++(0,0.45);
  \draw[primarydark,dashed,line width=0.5pt] (c2) -- ++(0,0.45);
  % amplitude
  \draw[<->,>=Stealth,formulacol,thick] (8.4,0) -- (8.4,\A)
        node[midway,right]{$A$};
\end{tikzpicture}

\vspace{0.6cm}

\begin{tcolorbox}[boxbase,colback=primary!5,colframe=primary,width=15cm,
    title={\textsc{Objectifs pédagogiques de la fiche}}]
À l'issue de ce cours, l'étudiant·e sera capable de :
\begin{itemize}[leftmargin=1.6em,topsep=2pt,itemsep=1pt]
  \item définir une \textbf{onde mécanique progressive} et distinguer ondes
        \textbf{longitudinale} et \textbf{transversale} ;
  \item maîtriser les \textbf{grandeurs caractéristiques} d'une onde
        (période, fréquence, longueur d'onde, célérité) et leurs relations ;
  \item écrire et exploiter l'\textbf{équation d'une onde sinusoïdale}
        (élongation, vitesse, accélération) et la notion de \textbf{déphasage} ;
  \item décrire les \textbf{ondes stationnaires}, leurs \textbf{nœuds} et
        \textbf{ventres}, et calculer les \textbf{harmoniques} d'une corde ;
  \item expliquer la \textbf{diffraction} et les \textbf{interférences}
        (franges, interfrange, conditions constructive/destructive).
\end{itemize}
\end{tcolorbox}

\vspace{0.45cm}
{\small\color{black!60} Document de synthèse rédigé d'après la Fiche 6 du livre — figures originales tracées en TikZ.}
\par
\endgroup

% ---------------------------------------------------------------------
%  TABLE DES MATIÈRES
% ---------------------------------------------------------------------
\newpage
\renewcommand{\contentsname}{\textcolor{primarydark}{Table des matières}}
\tableofcontents
\thispagestyle{fancy}
\newpage

% =====================================================================
\section{Propagation d'une onde}
% =====================================================================

\subsection{Notion d'onde mécanique progressive}

Dans la vie courante, on observe sans cesse des « signaux » qui voyagent : une vague
qui parcourt la surface de l'eau, le son d'une voix qui traverse une pièce, la secousse
qui file le long d'une corde à sauter. Dans tous ces cas, \emph{quelque chose se déplace
d'un point à un autre} — mais ce « quelque chose » n'est ni de la matière, ni un objet :
c'est une \textbf{perturbation} et l'\textbf{énergie} qu'elle transporte.

\begin{definition}[onde mécanique progressive]
Une \textbf{onde mécanique progressive} est la \emph{propagation d'une perturbation locale}
dans un \textbf{milieu élastique} (corde, ressort, eau, air, sol\ldots), \textbf{sans transport
de matière} mais \textbf{avec transport d'énergie}. Chaque point du milieu, lorsqu'il est
atteint, ne fait qu'osciller \emph{autour de sa position d'équilibre} : il ne « suit » pas l'onde,
il reproduit avec un certain retard le mouvement de la source.
\end{definition}

\begin{remarque}
Le mot \emph{élastique} est essentiel : le milieu doit pouvoir se déformer puis revenir à son
état initial. C'est cette élasticité qui transmet la perturbation de proche en proche. Le mot
\emph{progressive} signifie que la perturbation \emph{avance} (par opposition à l'onde
\emph{stationnaire} de la section~\ref{sec:stationnaire}, qui semble figée).
\end{remarque}

\subsection{Ondes longitudinales et ondes transversales}

Selon la \emph{direction de la perturbation par rapport à la direction de propagation}, on
distingue deux grandes familles d'ondes.

\begin{definition}[onde longitudinale]
Une onde est \textbf{longitudinale} lorsque les points du milieu vibrent
\emph{parallèlement} à la direction de propagation. L'onde se présente alors comme une
succession de zones de \textbf{compression} (points resserrés) et de \textbf{raréfaction}
(points écartés). Exemple type : la compression d'un ressort (figure~\ref{fig:ressort}), ou
le \textbf{son} dans l'air.
\end{definition}

\begin{figure}[h]
\centering
\begin{tikzpicture}[scale=1.0]
  % ---- ligne d'axe x ----
  \draw[->,>=Stealth,ondeA,line width=0.9pt] (-0.2,0) -- (10.6,0)
        node[right,black]{$x$ (m)};
  % ---- ligne verticale du point P ----
  \draw[dashed,black!70] (7.2,1.55) -- (7.2,-1.5) node[below]{};
  \node[font=\bfseries] at (7.2,1.9) {P};
  % ---- ressort du haut (instant t1) : compression à gauche ----
  \begin{scope}[yshift=1.05cm]
    \draw[black!75,line width=0.7pt,decorate,
          decoration={coil,aspect=0.6,segment length=2.0pt,amplitude=4.5pt}]
          (0,0) -- (1.7,0);
    \draw[black!75,line width=0.7pt,decorate,
          decoration={coil,aspect=0.6,segment length=5.0pt,amplitude=4.5pt}]
          (1.7,0) -- (10,0);
    \draw[->,>=Stealth,ondeA,line width=1.1pt] (0.35,0.42) -- (1.35,0.42)
          node[above,midway,black]{$\vv v$};
  \end{scope}
  % ---- ressort du bas (instant t2) : compression au niveau de P ----
  \begin{scope}[yshift=-1.05cm]
    \draw[black!75,line width=0.7pt,decorate,
          decoration={coil,aspect=0.6,segment length=5.0pt,amplitude=4.5pt}]
          (0,0) -- (6.4,0);
    \draw[black!75,line width=0.7pt,decorate,
          decoration={coil,aspect=0.6,segment length=2.0pt,amplitude=4.5pt}]
          (6.4,0) -- (8.0,0);
    \draw[black!75,line width=0.7pt,decorate,
          decoration={coil,aspect=0.6,segment length=5.0pt,amplitude=4.5pt}]
          (8.0,0) -- (10,0);
    \draw[->,>=Stealth,ondeA,line width=1.1pt] (5.9,0.42) -- (6.9,0.42)
          node[above,midway,black]{$\vv v$};
  \end{scope}
\end{tikzpicture}
\caption{L'onde se propage le long du ressort \textbf{parallèlement} à la direction de
propagation (onde longitudinale). La zone de compression se déplace vers la droite à la
vitesse $\vv v$ ; le point P, lui, ne fait qu'osciller sur place.}
\label{fig:ressort}
\end{figure}

\begin{definition}[onde transversale]
Une onde est \textbf{transversale} lorsque les points du milieu vibrent
\emph{perpendiculairement} à la direction de propagation. Exemple type : l'ébranlement
d'une corde tendue (figure~\ref{fig:corde}), ou les ondes à la surface de l'eau.
\end{definition}

\begin{figure}[h]
\centering
\begin{tikzpicture}[scale=1.0]
  % --- instant t1 : ébranlement à gauche, P au repos ---
  \begin{scope}[yshift=3.2cm]
    \draw[ondeA,line width=2pt] (0,-0.02) -- (10.2,-0.02);     % corde épaisse
    \draw[->,>=Stealth,black!70] (0,0) -- (0,1.25) node[left,black]{$y$};
    % impulsion (bosse) centrée en x=1
    \draw[ondeA,line width=2pt,smooth,samples=80,domain=0:2.4]
          plot (\x,{1.0*exp( -((\x-1)/0.42)^2 )});
    \node[ondeA] at (0.55,1.5) {Ébranlement};
    \draw[->,>=Stealth,ondeA,line width=1.1pt] (2.0,0.55) -- (3.0,0.55)
          node[above,midway,black]{$\vv v$};
    % graduations
    \foreach \i in {0,1,...,11}{\draw[black!60] (\i*0.9,0) -- (\i*0.9,-0.16);}
    \node[below] at (0,-0.2) {0}; \node[below] at (4.5,-0.2) {5};
    \node[below] at (9,-0.2) {10}; \node[right,black] at (10.2,0) {$x$ en m};
    \node[black] at (7.3,0.55) {Corde au repos};
    \draw[dashed,black!70] (4.5,1.35) -- (4.5,-0.05);
    \node[font=\bfseries] at (4.5,1.6) {P};
  \end{scope}
  % --- instant t2 : la bosse est arrivée en P ---
  \begin{scope}[yshift=0cm]
    \draw[ondeA,line width=2pt] (0,-0.02) -- (10.2,-0.02);
    \draw[->,>=Stealth,black!70] (0,0) -- (0,1.25) node[left,black]{$y$};
    \draw[ondeA,line width=2pt,smooth,samples=80,domain=3.1:5.9]
          plot (\x,{1.0*exp( -((\x-4.5)/0.42)^2 )});
    \draw[->,>=Stealth,ondeA,line width=1.1pt] (5.4,0.55) -- (6.4,0.55)
          node[above,midway,black]{$\vv v$};
    \foreach \i in {0,1,...,11}{\draw[black!60] (\i*0.9,0) -- (\i*0.9,-0.16);}
    \node[below] at (0,-0.2) {0}; \node[below] at (4.5,-0.2) {5};
    \node[below] at (9,-0.2) {10}; \node[right,black] at (10.2,0) {$x$ en m};
    \node[black] at (7.7,0.55) {Corde au repos};
    \draw[dashed,black!70] (4.5,1.35) -- (4.5,-0.05);
    \node[font=\bfseries] at (4.5,1.6) {P};
    \node[black] at (4.5,-0.95) {$v = x/t$};
  \end{scope}
\end{tikzpicture}
\caption{L'onde se propage le long de la corde élastique \textbf{perpendiculairement} à la
direction de propagation (onde transversale). Entre les deux instants, la bosse a parcouru
la distance jusqu'à P à la vitesse $v = x/t$.}
\label{fig:corde}
\end{figure}

\begin{theoreme}[indépendance de la célérité]
La \textbf{vitesse de propagation} (ou \textbf{célérité}) d'une onde est \emph{indépendante}
du type de perturbation (longitudinale ou transversale) et du \emph{sens} du déplacement.
Elle ne dépend que des \textbf{propriétés du milieu} (sa rigidité, sa masse linéique, sa
température\ldots). Une perturbation deux fois plus grande ne va donc pas plus vite : seule
son énergie change, pas sa célérité.
\end{theoreme}

\begin{exemple}[reconnaître le type d'onde]
\begin{itemize}[leftmargin=1.5em,topsep=2pt,itemsep=3pt]
  \item \textbf{Le son dans l'air} est une onde \emph{longitudinale} : l'air se comprime et
        se détend dans la \emph{même} direction que celle où le son voyage. C'est pour cela
        qu'on ne peut pas « polariser » le son.
  \item \textbf{La houle, la corde, la lumière} (onde électromagnétique) sont
        \emph{transversales} : la grandeur qui oscille (hauteur d'eau, position de la corde,
        champ électrique) est \emph{perpendiculaire} à la direction de propagation.
  \item \textbf{Le bateau au mouillage} qui monte et descend au passage d'une vague illustre
        parfaitement le « sans transport de matière » : le bateau revient à sa place, c'est
        seulement l'énergie de la vague qui a continué son chemin.
\end{itemize}
\end{exemple}

% =====================================================================
\section{Équation de propagation et grandeurs caractéristiques}
% =====================================================================

\subsection{Équation de propagation : le retard}

Considérons une source située à l'extrémité du milieu ($x=0$). Elle impose un mouvement
$y_0(t)$. Un point P situé plus loin, à l'abscisse $x$, va reproduire \emph{exactement} ce
mouvement, mais \textbf{en retard} : il faut que la perturbation parcoure la distance $x$ à la
célérité $v$, ce qui prend un temps $\tau = x/v$.

\begin{formule}[équation de propagation]
\[ \boxed{\,y_{\mathrm P}(t) = y_0\!\left(t - \dfrac{x}{v}\right)\,} \]
Le point P à l'abscisse $x$ refait, à l'instant $t$, ce que faisait la source à l'instant
antérieur $t - x/v$.
\begin{ou}
  \item $y_{\mathrm P}(t)$ : \textbf{élongation} (déplacement) du point P à l'instant $t$, en
        mètres (\si{\metre}) ;
  \item $y_0$ : loi horaire imposée par la \textbf{source} (en \si{\metre}) ;
  \item $x$ : \textbf{abscisse} (distance) du point P à la source, en mètres (\si{\metre}) ;
  \item $v$ : \textbf{célérité} de l'onde dans le milieu, en \si{\metre\per\second} ;
  \item $\dfrac{x}{v} = \tau$ : \textbf{retard} de P sur la source, en secondes (\si{\second}).
\end{ou}
\end{formule}

\noindent\textbf{Lecture.} Le terme $t - x/v$ est la signature même d'une onde qui se propage
\emph{vers les $x$ croissants}. Si l'onde se propageait dans l'autre sens, on écrirait
$t + x/v$. Tout point reçoit l'information « plus tard que la source », jamais avant : c'est le
principe de causalité.

\subsection{Grandeurs caractéristiques d'une onde périodique}

Lorsque la source répète indéfiniment le même mouvement, l'onde devient \textbf{périodique} :
elle se reproduit identique à elle-même à la fois \emph{dans le temps} (période $T$) et
\emph{dans l'espace} (longueur d'onde $\lambda$).

\begin{definition}[période, fréquence, longueur d'onde]
\begin{itemize}[leftmargin=1.5em,topsep=2pt,itemsep=2pt]
  \item La \textbf{période} $T$ est le temps nécessaire pour effectuer \emph{une vibration
        complète} (en \si{\second}).
  \item La \textbf{fréquence} $f$ est le \emph{nombre de vibrations par unité de temps}
        (en \si{\hertz}, $\SI{1}{\hertz} = \SI{1}{\per\second}$).
  \item La \textbf{longueur d'onde} $\lambda$ est la \emph{distance} entre deux points
        consécutifs de l'onde qui vibrent à l'identique (deux crêtes successives par
        exemple), en \si{\metre}. C'est la « période spatiale ».
\end{itemize}
\end{definition}

\begin{formule}[période et fréquence]
\[ \boxed{\,T = \dfrac{1}{f}\,}\qquad\Longleftrightarrow\qquad f = \dfrac{1}{T} \]
\begin{ou}
  \item $T$ : période, en secondes (\si{\second}) ;
  \item $f$ : fréquence, en hertz (\si{\hertz}) ou \si{\per\second}.
\end{ou}
Période et fréquence sont \textbf{inverses} l'une de l'autre : plus une onde vibre vite
(grande fréquence), plus sa période est courte.
\end{formule}

\begin{formule}[longueur d'onde et célérité]
\[ \boxed{\,\lambda = v\times T = \dfrac{v}{f}\,}
   \qquad\text{et de façon équivalente}\qquad
   \boxed{\,v = \lambda\times f\,} \]
\begin{ou}
  \item $\lambda$ : longueur d'onde, en mètres (\si{\metre}) ;
  \item $v$ : célérité de l'onde, en \si{\metre\per\second} ;
  \item $T$ : période, en secondes (\si{\second}) ; \quad $f$ : fréquence, en \si{\hertz}.
\end{ou}
\textbf{Interprétation.} Pendant une période $T$, la source fait une oscillation complète et,
dans le même temps, l'onde avance d'exactement une longueur d'onde $\lambda$. La relation
$\lambda = vT$ n'est donc que la définition « distance = vitesse $\times$ temps » appliquée sur
une période.
\end{formule}

\begin{attention}
$f$ (la fréquence) est imposée par la \textbf{source} et ne change pas quand l'onde change de
milieu. En revanche $v$ et $\lambda$ \emph{dépendent du milieu}. Quand une onde passe d'un
milieu à un autre (l'air vers l'eau par ex.), sa fréquence reste la même mais sa longueur
d'onde change, puisque $\lambda = v/f$.
\end{attention}

\begin{exemple}[calcul des grandeurs d'une onde sonore]
Un haut-parleur émet un son pur de fréquence $f = \SI{100}{\hertz}$ dans l'air, où la célérité
du son vaut $v = \SI{340}{\metre\per\second}$. Déterminons toutes les grandeurs.

\smallskip
\textbf{Période.} \quad
$\displaystyle T = \frac{1}{f} = \frac{1}{\SI{100}{\hertz}} = \SI{0.01}{\second} = \SI{10}{\milli\second}.$
Chaque oscillation complète de l'air dure un centième de seconde.

\smallskip
\textbf{Longueur d'onde.} \quad
$\displaystyle \lambda = \frac{v}{f} = \frac{\SI{340}{\metre\per\second}}{\SI{100}{\hertz}}
 = \SI{3.4}{\metre}.$
Deux zones de compression successives sont distantes de \SI{3.4}{\metre} dans l'air.

\smallskip
\textbf{Vérification.} \quad
$v = \lambda f = \SI{3.4}{\metre}\times\SI{100}{\hertz} = \SI{340}{\metre\per\second}$ : on
retrouve bien la célérité de départ. \\[2pt]
\textbf{Et dans l'eau ?} La célérité du son y vaut environ \SI{1480}{\metre\per\second}. La
fréquence reste $\SI{100}{\hertz}$ (c'est la source qui la fixe), donc
$\lambda_{\text{eau}} = 1480/100 = \SI{14.8}{\metre}$ : la longueur d'onde est plus de
quatre fois plus grande, alors que le son « est le même » pour la source.
\end{exemple}

\begin{exemple}[onde à la surface de l'eau]
Sur un lac, on compte que les crêtes des vagues passent toutes les $T = \SI{4}{\second}$
sous un ponton, et l'on mesure une distance de $\lambda = \SI{2}{\metre}$ entre deux crêtes
voisines. Alors :
\[ v = \frac{\lambda}{T} = \frac{\SI{2}{\metre}}{\SI{4}{\second}} = \SI{0.5}{\metre\per\second},
   \qquad f = \frac{1}{T} = \frac{1}{\SI{4}{\second}} = \SI{0.25}{\hertz}. \]
Les vagues avancent donc à un demi-mètre par seconde, en effectuant un quart d'oscillation
par seconde. On note au passage qu'une \emph{grande} longueur d'onde associée à une
\emph{petite} fréquence donne une onde lente : tout est cohérent avec $v = \lambda f$.
\end{exemple}

% =====================================================================
\section{Caractéristiques d'une onde sinusoïdale}
% =====================================================================
\label{sec:sinus}

Le cas le plus important est celui où la source effectue un mouvement \textbf{sinusoïdal}
(harmonique simple). L'onde qui en résulte est alors décrite par une fonction sinus, à la fois
du temps $t$ et de la position $x$.

\subsection{Équation de l'élongation}

\begin{formule}[élongation d'une onde sinusoïdale progressive]
\[ \boxed{\,y_{\mathrm P}(t) = A\,\sin\!\left[\omega\!\left(t-\frac{x}{v}\right)\right]
   = A\,\sin\!\left(2\pi f\,t - \frac{2\pi x}{Tv}\right)
   = A\,\sin\!\left(\omega t - k x\right)\,} \]
\begin{ou}
  \item $y_{\mathrm P}(t)$ : élongation du point P à l'instant $t$, en \si{\metre} ;
  \item $A$ : \textbf{amplitude}, élongation maximale (toujours positive), en \si{\metre} ;
  \item $\omega$ : \textbf{pulsation} (ou fréquence angulaire), en \si{\radian\per\second} ;
  \item $k$ : \textbf{nombre d'onde} (pulsation spatiale), en \si{\radian\per\metre} ;
  \item $t$ : temps, en \si{\second} ; \quad $x$ : abscisse du point, en \si{\metre} ;
  \item $\omega t - kx$ : \textbf{phase} de l'onde, en \si{\radian}.
\end{ou}
\end{formule}

\begin{formule}[pulsation et nombre d'onde]
\[ \boxed{\,\omega = 2\pi f = \frac{2\pi}{T}\,}\qquad\qquad
   \boxed{\,k = \frac{2\pi}{\lambda}\,} \]
\begin{ou}
  \item $\omega$ : pulsation, en \si{\radian\per\second} — c'est le « nombre de radians
        balayés par seconde ». Un tour complet ($2\pi$ rad) correspond à une période ;
  \item $k$ : nombre d'onde, en \si{\radian\per\metre} — le « nombre de radians par mètre ».
        Un tour complet correspond à une longueur d'onde ;
  \item on en déduit aussi $\dfrac{\omega}{k} = \dfrac{2\pi f}{2\pi/\lambda} = \lambda f = v$ :
        \textbf{la célérité est le rapport de la pulsation au nombre d'onde}.
\end{ou}
\end{formule}

La figure~\ref{fig:sinus} représente l'onde « photographiée » à un instant donné en fonction
de la position $x$. On y lit l'amplitude $A$ (hauteur des crêtes) et la longueur d'onde
$\lambda$ (distance entre deux crêtes).

\begin{figure}[h]
\centering
\begin{tikzpicture}[scale=1.0]
  \def\A{1.5}\def\Lc{3.0}   % amplitude (cm), longueur d'onde (cm)
  % axes
  \draw[->,>=Stealth,black!75] (-0.3,0) -- (10.2,0) node[right]{$x$ (m)};
  \draw[->,>=Stealth,black!75] (0,-2.0) -- (0,2.1) node[left]{$y$};
  \draw[dashed,black!40] (-0.2,\A) -- (10,\A);
  \draw[dashed,black!40] (-0.2,-\A) -- (10,-\A);
  \node[left] at (0,\A) {$A$}; \node[left] at (0,-\A) {$-A$};
  \node[below left] at (0,0) {$0$};
  % courbe sinusoïdale : crêtes en x = L/4 + nL
  \draw[primary,line width=1.5pt,smooth,samples=240,domain=0:9.6]
        plot (\x,{\A*sin( \x*(360/\Lc) )});
  % points remarquables
  \coordinate (A) at ({0.25*\Lc},\A);    % crête 1
  \coordinate (C) at ({1.25*\Lc},\A);    % crête 2  (concordance avec A)
  \coordinate (B) at ({0.75*\Lc},-\A);   % creux 1  (opposition avec A)
  \coordinate (D) at ({1.75*\Lc},-\A);   % creux 2
  \coordinate (P) at ({1.55*\Lc},{\A*sin(1.55*360)});
  \foreach \pt/\lab/\pos in {A/A/above,C/C/above,B/B/below,D/D/below}{
     \fill[primarydark] (\pt) circle (2.3pt); \node[\pos] at (\pt) {\textbf{\lab}};}
  \fill[formulacol] (P) circle (2.3pt); \node[right,formulacol] at (P) {\textbf{P}};
  % longueur d'onde entre A et C
  \draw[<->,>=Stealth,primarydark,thick] ($(A)+(0,0.4)$) -- ($(C)+(0,0.4)$)
       node[midway,above,fill=white,inner sep=1pt]{$\lambda$};
  \draw[primarydark,dashed,line width=0.5pt] (A)--($(A)+(0,0.5)$);
  \draw[primarydark,dashed,line width=0.5pt] (C)--($(C)+(0,0.5)$);
  % demi-longueur d'onde entre A et B (opposition)
  \draw[<->,>=Stealth,formulacol] ($(A)+(0,-0.0)$)
        .. controls ($(A)!0.5!(B)+(0,1.0)$) .. ($(B)+(0,0.0)$);
  \node[formulacol,fill=white,inner sep=1pt] at ($(A)!0.5!(B)+(0,0.55)$) {$\lambda/2$};
\end{tikzpicture}
\caption{Onde sinusoïdale à un instant fixé. Les points \textbf{A} et \textbf{C}, distants de
$\lambda$, vibrent \emph{en concordance de phase}. Les points \textbf{A} et \textbf{B},
distants de $\lambda/2$, vibrent \emph{en opposition de phase}.}
\label{fig:sinus}
\end{figure}

\subsection{Vitesse et accélération d'un point du milieu}

Un point P du milieu oscille : à tout instant il possède une \textbf{vitesse} et une
\textbf{accélération}, obtenues en dérivant l'élongation par rapport au temps. \emph{Attention}
à ne pas confondre cette vitesse $v_{\mathrm P}$ (mouvement \emph{transversal} du point) avec la
célérité $v$ (vitesse de \emph{propagation} de l'onde).

\begin{formule}[vitesse d'un point du milieu]
\[ \boxed{\,v_{\mathrm P}(t) = y_{\mathrm P}'(t)
   = A\,\omega\,\cos\!\left[\omega\!\left(t-\frac{x}{v}\right)\right]\,} \]
\begin{ou}
  \item $v_{\mathrm P}(t)$ : vitesse d'oscillation du point P, en \si{\metre\per\second} ;
  \item $A\,\omega$ : \textbf{vitesse maximale} du point (atteinte quand il passe par sa
        position d'équilibre), en \si{\metre\per\second} ;
  \item $A$ en \si{\metre}, $\omega$ en \si{\radian\per\second}.
\end{ou}
\end{formule}

\begin{formule}[accélération d'un point du milieu]
\[ \boxed{\,a_{\mathrm P}(t) = v_{\mathrm P}'(t)
   = -A\,\omega^{2}\,\sin\!\left[\omega\!\left(t-\frac{x}{v}\right)\right]
   = -\,\omega^{2}\,y_{\mathrm P}(t)\,} \]
\begin{ou}
  \item $a_{\mathrm P}(t)$ : accélération du point P, en \si{\metre\per\second\squared} ;
  \item $A\,\omega^{2}$ : \textbf{accélération maximale} (atteinte aux crêtes et aux creux), en
        \si{\metre\per\second\squared} ;
  \item la relation $a_{\mathrm P} = -\omega^{2}y_{\mathrm P}$ est la signature du
        \textbf{mouvement harmonique simple} : l'accélération est toujours \emph{opposée} à
        l'élongation et proportionnelle à elle (force de rappel).
\end{ou}
\end{formule}

\subsection{Déphasage, concordance et opposition de phase}

Si l'onde est \textbf{déphasée} d'un angle $\varphi$ (par exemple parce que la source n'a pas
démarré à $t=0$), son équation devient :
\[ \boxed{\,y_{\mathrm P}(t) = A\,\sin(\omega t - k x + \varphi)\,} \]
où $\varphi$ est la \textbf{phase à l'origine}, en radians (\si{\radian}).

\begin{definition}[concordance et opposition de phase]
On compare deux points (ou deux ondes) selon la distance $d$ qui les sépare :
\begin{itemize}[leftmargin=1.5em,topsep=2pt,itemsep=2pt]
  \item \textbf{Concordance de phase} : si $d = n\lambda$ ($n\in\mathbb{Z}$), les deux points
        vibrent \emph{exactement de la même façon} au même instant (ex. : A et C
        sur la figure~\ref{fig:sinus}).
  \item \textbf{Opposition de phase} : si $d = (2n+1)\dfrac{\lambda}{2}$ ($n\in\mathbb{Z}$),
        les deux points vibrent \emph{de manière inverse} : quand l'un monte, l'autre descend
        (ex. : A et B).
\end{itemize}
\end{definition}

\begin{exemple}[déphasage entre deux ondes — étude complète]
On considère deux ondes 1 et 2 de \emph{même fréquence} se propageant de gauche à droite
(figure~\ref{fig:deux-ondes}). L'\textbf{Onde 1} a pour équation $y_1 = 0{,}2\,\sin(2\pi t)$.

\begin{center}
\begin{tikzpicture}
\begin{axis}[
  width=13.4cm,height=6.8cm,
  xmin=0,xmax=3,ymin=-0.37,ymax=0.37,
  xtick={0,0.5,1,1.5,2,2.5,3},
  ytick={-0.3,-0.2,-0.1,0,0.1,0.2,0.3},
  grid=both,grid style={black!12},
  axis lines=left,
  xlabel={Temps en s},ylabel={Amplitude en m},
  xlabel style={at={(axis description cs:0.5,-0.09)},anchor=north},
  ylabel style={at={(axis description cs:-0.065,0.5)},rotate=0,anchor=south,
     sloped=false},
  tick label style={font=\footnotesize},
  label style={font=\footnotesize},
  samples=200,clip=false,
  legend cell align=left,
  legend style={at={(0.5,1.07)},anchor=south,legend columns=2,
     draw=black!30,font=\footnotesize\bfseries,column sep=10pt},
]
  \addplot[ondeA,line width=1.3pt,domain=0:3]{0.2*sin(deg(2*pi*x))};
  \addlegendentry{Onde 1}
  \addplot[ondeB,line width=1.3pt,domain=0:3]{0.2*sin(deg(2*pi*x)-90)};
  \addlegendentry{Onde 2}
  % ligne de zéro plus marquée
  \draw[black!45,line width=0.6pt] (axis cs:0,0) -- (axis cs:3,0);
  % écart de phase Delta t entre les deux premières crêtes (0,25 -> 0,50)
  \draw[<->,>=Stealth,black] (axis cs:0.25,0.265) -- (axis cs:0.5,0.265);
  \node[font=\footnotesize,anchor=south,inner sep=1.5pt]
       at (axis cs:0.37,0.275) {$\Delta t = \SI{0.25}{\second}$};
  % période T entre deux crêtes successives de l'Onde 1 (0,25 -> 1,25)
  \draw[<->,>=Stealth,black] (axis cs:0.25,-0.275) -- (axis cs:1.25,-0.275);
  \node[font=\footnotesize,anchor=north,inner sep=1.5pt]
       at (axis cs:0.75,-0.285) {$T = \SI{1}{\second}$};
\end{axis}
\end{tikzpicture}
\captionof{figure}{Les deux ondes ont la même période $T=\SI{1}{\second}$ et la même amplitude
$A=\SI{0.2}{\metre}$. L'Onde 2 est décalée de $\Delta t = \SI{0.25}{\second}$ vers la droite :
elle est \emph{en retard} sur l'Onde 1.}
\label{fig:deux-ondes}
\end{center}

\textbf{1. Identification de l'avance/retard.} L'Onde 1 atteint son maximum
($A = \SI{0.2}{\metre}$) \emph{en premier} : elle est donc \textbf{en avance de phase} sur
l'Onde 2. L'écart temporel entre les deux (par exemple entre leurs premiers ventres) est
$\Delta t = \SI{0.25}{\second}$. L'Onde 2 étant en retard, son déphasage $\varphi_2$ sera
\emph{négatif} : $\Delta\varphi = \varphi_2 - \varphi_1 < 0$.

\smallskip
\textbf{2. Calcul du déphasage.} On utilise une simple proportionnalité : une période entière
$T$ correspond à un tour complet $2\pi$. Donc un retard $\Delta t$ correspond à un angle :
\[
\begin{array}{r@{\;\longrightarrow\;}l}
 T = \SI{1}{\second} & 2\pi \\[2pt]
 \Delta t = \SI{0.25}{\second} & |\Delta\varphi|\;?
\end{array}
\qquad\Longrightarrow\qquad
|\Delta\varphi| = \frac{\Delta t\times 2\pi}{T}
 = \frac{0{,}25\times 2\pi}{1} = \frac{\pi}{2}.
\]
Comme l'Onde 2 est en \emph{retard}, on prend le signe négatif :
\[ \Delta\varphi = \varphi_2 - \varphi_1 = \varphi_2 - 0
   = -\left|\frac{\Delta t\times 2\pi}{T}\right| = -\frac{\pi}{2}
   \quad\Longrightarrow\quad \varphi_2 = -\frac{\pi}{2}. \]

\textbf{3. Équation de l'Onde 2.} On reporte $\varphi_2$ dans l'équation générale (avec la
même pulsation $\omega = 2\pi$) :
\[ \boxed{\,y_2 = 0{,}2\,\sin\!\left(2\pi t + \varphi_2\right)
   = 0{,}2\,\sin\!\left(2\pi t - \frac{\pi}{2}\right).} \]

\textbf{Important.} Si, au contraire, l'Onde 2 avait été \emph{en avance} sur l'Onde 1, on
aurait eu $\Delta\varphi = \varphi_2 - \varphi_1 > 0$ (déphasage positif).
\end{exemple}

% =====================================================================
\section{Ondes stationnaires}
% =====================================================================
\label{sec:stationnaire}

\subsection{Réflexion et formation d'une onde stationnaire}

\begin{definition}[réflexion sur un obstacle fixe]
Lorsqu'une onde progressive arrive sur un \textbf{obstacle fixe}, dans un milieu élastique
inchangé, elle subit une \textbf{réflexion totale} \emph{avec changement de signe de
l'élongation} (la bosse repart en creux). L'énergie est \textbf{intégralement conservée} :
l'onde réfléchie repart à la \emph{même célérité} que l'onde incidente.
\end{definition}

\begin{figure}[h]
\centering
\begin{tikzpicture}[scale=1.0]
  % --- onde incidente (haut) ---
  \begin{scope}[yshift=1.6cm]
    \draw[black,line width=1.6pt] (0,0) -- (10.0,0);
    \draw[black,line width=1.6pt] (0,-0.45) -- (0,0.9);          % mur gauche
    \draw[ondeA,line width=1.8pt,smooth,samples=70,domain=0.1:1.9]
          plot (\x,{0.9*exp(-((\x-1)/0.4)^2)});
    \node[ondeA,font=\bfseries\footnotesize] at (1.0,1.25) {Onde incidente};
    \draw[->,>=Stealth,ondeA,line width=1.1pt] (1.7,0.45)--(2.7,0.45)
          node[above,midway,black]{$\vv v$};
    \foreach \i in {0,...,10}{\draw[black!60] (\i,0)--(\i,-0.18);}
    \node[below] at (0,-0.2){0}; \node[below] at (5,-0.2){5};
    \node[below] at (10,-0.2){10}; \node[right] at (10.0,0){$x$ en m};
    \node at (7.4,0.45){Corde au repos};
  \end{scope}
  % --- onde réfléchie (bas) : revient en creux, vers la gauche ---
  \begin{scope}[yshift=-1.1cm]
    \draw[black,line width=1.6pt] (0,0) -- (10.0,0);
    \draw[black,line width=1.6pt] (10,-0.9) -- (10,0.45);        % mur droit (obstacle)
    \draw[ondeB,line width=1.8pt,smooth,samples=70,domain=8.1:9.9]
          plot (\x,{-0.9*exp(-((\x-9)/0.4)^2)});
    \node[ondeB,font=\bfseries\footnotesize] at (8.7,-1.25) {Onde réfléchie};
    \draw[->,>=Stealth,ondeB,line width=1.1pt] (8.3,-0.45)--(7.3,-0.45)
          node[above,midway,black]{$\vv v$};
    \foreach \i in {0,...,10}{\draw[black!60] (\i,0)--(\i,0.18);}
    \node[above] at (0,0.2){0}; \node[above] at (5,0.2){5};
    \node[above] at (10,0.2){10}; \node[right] at (10.0,0){$x$ en m};
    \node at (2.6,0.45){Corde au repos};
  \end{scope}
\end{tikzpicture}
\caption{Réflexion d'un ébranlement sur l'obstacle fixe (mur de droite) : la bosse incidente
(bleue) repart en creux (orange), avec la même vitesse. L'énergie est conservée.}
\label{fig:reflexion}
\end{figure}

\begin{definition}[onde stationnaire, nœuds et ventres]
La \textbf{superposition} (interférence) de l'onde incidente et de l'onde réfléchie donne
naissance à une \textbf{onde stationnaire} : une figure qui ne se déplace plus, mais
\emph{oscille sur place}. Elle est caractérisée par :
\begin{itemize}[leftmargin=1.5em,topsep=2pt,itemsep=2pt]
  \item des \textbf{nœuds} : points d'\emph{élongation toujours nulle} (le milieu y reste
        immobile) ;
  \item des \textbf{ventres} : points d'\emph{élongation maximale} (le milieu y oscille avec
        l'amplitude la plus grande).
\end{itemize}
\end{definition}

\begin{figure}[h]
\centering
\begin{tikzpicture}[scale=1.0]
  \def\A{1.1}\def\Ls{9.0}
  % axes
  \draw[->,>=Stealth,black!75] (-0.2,0) -- (10.4,0) node[right]{};
  \draw[->,>=Stealth,black!75] (0,-1.95) -- (0,2.0) node[left]{Amplitude};
  % obstacle fixe (mur droit)
  \draw[black,line width=1.4pt] (\Ls,-1.35) -- (\Ls,1.35);
  \node[anchor=west,font=\footnotesize] at (\Ls+0.1,1.1) {obstacle fixe};
  % ondes incidente (bleu) et réfléchie (rouge), faibles
  \draw[ondeA!60,line width=1pt,smooth,samples=240,domain=0:\Ls]
        plot (\x,{0.58*sin(\x*60)});
  \draw[ondeB!60,line width=1pt,smooth,samples=240,domain=0:\Ls]
        plot (\x,{-0.58*sin(\x*60)});
  % onde stationnaire (noir épais)
  \draw[black,line width=1.8pt,smooth,samples=260,domain=0:\Ls]
        plot (\x,{\A*sin(\x*60)});
  % noeuds (points) et flèches de ventre
  \foreach \xn in {0,3,6,9}{\fill[primarydark] (\xn,0) circle (2.3pt);}
  \foreach \xv/\dir in {1.5/1,4.5/-1,7.5/1}{
     \draw[->,>=Stealth,black!70,line width=0.7pt]
        (\xv,{0.15*\dir}) -- (\xv,{(\A-0.05)*\dir});}
  % --- étiquettes dans les marges (sans recouvrir la courbe noire) ---
  % ventres (haut) en x = 1,5 et 7,5
  \node[font=\footnotesize] (V1) at (1.5,1.66) {Ventre};
  \draw[black!55,line width=0.5pt] (V1) -- (1.5,1.2);
  \node[font=\footnotesize] (V3) at (7.5,1.66) {Ventre};
  \draw[black!55,line width=0.5pt] (V3) -- (7.5,1.2);
  % noeuds (haut) en x = 3 et 6
  \node[font=\footnotesize] (N1) at (3,1.66) {Nœud};
  \draw[black!55,line width=0.5pt] (N1) -- (3,0.12);
  \node[font=\footnotesize] (N2) at (6,1.66) {Nœud};
  \draw[black!55,line width=0.5pt] (N2) -- (6,0.12);
  % ventre (bas) en x = 4,5
  \node[font=\footnotesize] (V2) at (4.5,-1.66) {Ventre};
  \draw[black!55,line width=0.5pt] (V2) -- (4.5,-1.2);
  % identification des courbes (codée par la couleur)
  \node[ondeA!75!black,font=\footnotesize] at (1.95,-1.66) {Onde incidente};
  \node[ondeB!75!black,font=\footnotesize] at (7.15,-1.66) {Onde réfléchie};
\end{tikzpicture}
\caption{Onde stationnaire (trait noir épais) résultant de la somme de l'onde incidente
(bleue) et de l'onde réfléchie (rouge) \emph{en phase}. Les \textbf{nœuds} sont fixes
(amplitude nulle), les \textbf{ventres} oscillent au maximum.}
\label{fig:stationnaire}
\end{figure}

\begin{remarque}
\begin{itemize}[leftmargin=1.5em,topsep=2pt,itemsep=2pt]
  \item L'amplitude de l'onde stationnaire est simplement la \textbf{somme des amplitudes} de
        chacune des deux ondes (incidente et réfléchie).
  \item On dit que l'onde stationnaire est \textbf{résonante} lorsque les conditions de
        longueur d'onde (ou de fréquence) sont fixées par la géométrie du milieu.
\end{itemize}
\end{remarque}

\subsection{Harmoniques d'une corde de longueur \texorpdfstring{$L$}{L}}

Pour une corde de longueur $L$ fixée, seules certaines longueurs d'onde « rentrent » exactement
dans la corde : ce sont les \textbf{modes propres} ou \textbf{harmoniques}.

\begin{formule}[corde fixée aux deux extrémités]
\[ \boxed{\,\lambda_n = \frac{2L}{n}\,}\qquad
   \boxed{\,f_n = \frac{c}{\lambda_n} = \frac{n\,c}{2L}\,}\qquad (n = 1,2,3,\dots) \]
\begin{ou}
  \item $\lambda_n$ : longueur d'onde du mode $n$, en \si{\metre} ;
  \item $L$ : longueur de la corde, en \si{\metre} ;
  \item $c$ : célérité de l'onde sur la corde, en \si{\metre\per\second} ;
  \item $f_n$ : fréquence du mode $n$ (la $n$-ième harmonique), en \si{\hertz} ;
  \item $n$ : numéro de l'harmonique (entier $\ge 1$). $n=1$ : \textbf{mode fondamental}
        (1\textsuperscript{re} harmonique) ; $n=2$ : 2\textsuperscript{e} harmonique ; etc.
\end{ou}
Les deux extrémités fixes \emph{imposent un nœud à chaque bout} : la corde contient donc un
nombre entier de « fuseaux » (demi-longueurs d'onde), d'où $L = n\,\lambda_n/2$.
\end{formule}

\begin{figure}[h]
\centering
\begin{tikzpicture}[scale=1.0]
  \def\Lh{9.0}\def\Ah{0.62}
  \foreach \n/\yy/\lab in {1/3.6/{$n=1$}, 2/1.8/{$n=2$}, 3/0.0/{$n=3$}}{
    \begin{scope}[yshift=\yy cm]
      % murs (extrémités fixes) en vert
      \draw[excol,line width=3pt] (0,-0.75)--(0,0.75);
      \draw[excol,line width=3pt] (\Lh,-0.75)--(\Lh,0.75);
      % enveloppe supérieure (bleue) et inférieure (noire)
      \draw[defcol,line width=1.1pt,smooth,samples=200,domain=0:\Lh]
            plot (\x,{\Ah*sin(\x*(\n*180/\Lh))});
      \draw[black,line width=1.1pt,smooth,samples=200,domain=0:\Lh]
            plot (\x,{-\Ah*sin(\x*(\n*180/\Lh))});
      % noeuds
      \foreach \k in {0,...,\n}{\fill[black] (\k*\Lh/\n,0) circle (2.4pt);}
      \node[left=6pt,font=\small] at (0,0) {\lab};
    \end{scope}
  }
  % cote L sous la première portée
  \draw[<->,>=Stealth,black!75,dashed] (0,2.95) -- (\Lh,2.95)
        node[midway,fill=white,inner sep=1pt]{$L$};
\end{tikzpicture}
\caption{Trois premiers modes propres d'une corde fixée à ses deux extrémités. Le mode $n$
possède $n$ fuseaux (ventres) et $n+1$ nœuds (points noirs), dont les deux extrémités.}
\label{fig:harmoniques}
\end{figure}

\begin{formule}[corde avec une extrémité libre]
\[ \boxed{\,\lambda_n = \frac{4L}{2n-1}\,}\qquad (n = 1,2,3,\dots) \]
\begin{ou}
  \item $\lambda_n$ : longueur d'onde du mode $n$, en \si{\metre} ;
  \item $L$ : longueur de la corde, en \si{\metre} ;
  \item $n$ : numéro de l'harmonique (entier $\ge 1$).
\end{ou}
Une extrémité libre impose un \emph{ventre} (et non un nœud) à ce bout : seuls les modes
« impairs » subsistent, d'où le facteur $2n-1$ au dénominateur.
\end{formule}

\begin{exemple}[harmoniques d'une corde de guitare]
Une corde de longueur $L = \SI{1.2}{\metre}$, fixée à ses deux extrémités, est le siège d'ondes
de célérité $c = \SI{120}{\metre\per\second}$. Calculons les trois premiers harmoniques avec
$\lambda_n = 2L/n$ et $f_n = nc/(2L)$.

\smallskip
\textbf{Fondamental} ($n=1$) :
\[ \lambda_1 = \frac{2\times 1{,}2}{1} = \SI{2.4}{\metre},\qquad
   f_1 = \frac{1\times 120}{2\times 1{,}2} = \frac{120}{2{,}4} = \SI{50}{\hertz}. \]

\textbf{2\textsuperscript{e} harmonique} ($n=2$) :
\[ \lambda_2 = \frac{2\times 1{,}2}{2} = \SI{1.2}{\metre},\qquad
   f_2 = \frac{2\times 120}{2{,}4} = \SI{100}{\hertz} = 2 f_1. \]

\textbf{3\textsuperscript{e} harmonique} ($n=3$) :
\[ \lambda_3 = \frac{2\times 1{,}2}{3} = \SI{0.8}{\metre},\qquad
   f_3 = \frac{3\times 120}{2{,}4} = \SI{150}{\hertz} = 3 f_1. \]

\smallskip
\textbf{Conclusion.} Les fréquences propres sont des \emph{multiples entiers} du fondamental :
$f_n = n f_1 = 50\,n$ \si{\hertz}. C'est exactement ce qui donne aux instruments à cordes leur
sonorité « musicale » : la note de hauteur $f_1$ est accompagnée de ses harmoniques
$2f_1, 3f_1,\dots$ Plus la corde est courte (petit $L$) ou tendue (grand $c$), plus le
fondamental est aigu.
\end{exemple}

% =====================================================================
\section{Diffraction et interférences}
% =====================================================================

\subsection{Diffraction}

\begin{definition}[diffraction]
La \textbf{diffraction} est le \emph{changement de direction} d'une onde qui rencontre un
\textbf{obstacle} ou une \textbf{ouverture}. Elle n'est nettement observable que si la
dimension de l'obstacle (ou de la fente) est du \emph{même ordre de grandeur} que la longueur
d'onde $\lambda$. Après l'obstacle, l'onde « contourne » et s'étale : des fronts d'onde plans
deviennent des fronts d'onde quasi circulaires.
\end{definition}

\begin{figure}[h]
\centering
\begin{tikzpicture}[scale=1.0]
  % --- fronts d'onde plans (à gauche) ---
  \foreach \x in {0,0.6,1.2,1.8,2.4}{
     \draw[black,line width=1pt] (\x,-2)--(\x,2);}
  \draw[<->,>=Stealth,black] (1.2,2.25)--(1.8,2.25)
        node[midway,above,font=\footnotesize]{$\lambda$};
  \draw[->,>=Stealth,black,line width=1pt] (1.0,-2.45)--(2.2,-2.45)
        node[right,font=\footnotesize]{$\vv v$};
  \node[font=\footnotesize] at (1.2,-2.85) {Propagation des ondes};
  % --- barrière avec fente au centre ---
  \draw[black,line width=2.4pt] (3.4,0.45)--(3.4,2.1);
  \draw[black,line width=2.4pt] (3.4,-2.1)--(3.4,-0.45);
  \node[font=\footnotesize,rotate=90] at (3.05,1.25) {Obstacle};
  \node[font=\footnotesize,rotate=90] at (3.05,-1.25) {Obstacle};
  % --- fronts circulaires (à droite), centrés sur la fente ---
  \foreach \r in {0.55,1.1,1.65,2.2,2.75}{
     \draw[black,dashed,line width=0.8pt]
       ([shift={(3.4,0)}]75:\r) arc (75:-75:\r);}
  \draw[<->,>=Stealth,black] ([shift={(3.4,0)}]35:1.1)--([shift={(3.4,0)}]35:1.65)
        node[midway,above right,font=\footnotesize]{$\lambda$};
\end{tikzpicture}
\caption{Diffraction par une fente de taille comparable à $\lambda$ : les fronts d'onde plans
(à gauche) deviennent des fronts d'onde circulaires (à droite), comme si la fente était une
nouvelle source.}
\label{fig:diffraction}
\end{figure}

\subsection{Interférences de deux sources}

\begin{definition}[interférences]
Quand deux ondes de \emph{même longueur d'onde} issues de deux \textbf{sources synchrones}
(cohérentes) se \textbf{superposent}, leurs élongations s'ajoutent. Selon les endroits, elles
se renforcent ou s'annulent : on observe des \textbf{franges d'interférence}.
\begin{itemize}[leftmargin=1.5em,topsep=2pt,itemsep=2pt]
  \item \textbf{Interférence constructive} : les ondes arrivent \emph{en concordance de phase},
        elles se renforcent $\Rightarrow$ frange brillante / \textbf{ventre}.
  \item \textbf{Interférence destructive} : les ondes arrivent \emph{en opposition de phase},
        elles s'annulent mutuellement $\Rightarrow$ frange sombre / \textbf{nœud}.
\end{itemize}
\end{definition}

On considère deux sources synchrones $S_1$ et $S_2$ séparées d'une distance $a$, et un point
$M$ d'un écran placé à la distance $D$ (figure~\ref{fig:young}). On note $d_1 = S_1M$ et
$d_2 = S_2M$. L'élongation au point $M$ est la \emph{somme} des deux élongations :
\begin{formule}[superposition au point M]
\[ y_M = y_M^{1} + y_M^{2},\qquad
   \begin{cases}
     y_M^{1} = A\,\sin\!\Big(2\pi\big(t f - d_1/\lambda\big)\Big)\\[4pt]
     y_M^{2} = A\,\sin\!\Big(2\pi\big(t f - d_2/\lambda\big)\Big)
   \end{cases} \]
\begin{ou}
  \item $y_M^{1},\,y_M^{2}$ : élongations dues à chaque source au point $M$, en \si{\metre} ;
  \item $A$ : amplitude commune des deux sources, en \si{\metre} ;
  \item $d_1,\,d_2$ : distances des sources $S_1$, $S_2$ au point $M$, en \si{\metre} ;
  \item $f$ : fréquence commune (\si{\hertz}) ; $\lambda$ : longueur d'onde (\si{\metre}).
\end{ou}
\end{formule}

C'est la \textbf{différence de marche} $\delta = d_1 - d_2$ qui décide du type d'interférence :

\begin{theoreme}[conditions d'interférence]
\[ \text{\bf Constructive (ventre) : } d_1 - d_2 = n\lambda
   \qquad\qquad
   \text{\bf Destructive (nœud) : } d_1 - d_2 = (2n+1)\frac{\lambda}{2}, \]
avec $n\in\mathbb{Z}$. Autrement dit : si la différence de marche est un \emph{multiple entier}
de $\lambda$, les ondes se renforcent ; si c'est un \emph{multiple impair} de $\lambda/2$,
elles s'annulent.
\end{theoreme}

\begin{figure}[h]
\centering
\begin{tikzpicture}[scale=1.0]
  % plan des sources (gauche) et écran (droite)
  \draw[black,line width=1pt] (0,-2.6)--(0,2.8);
  \draw[black,line width=1.4pt] (7,-2.8)--(7,3.2);
  \node[right,font=\footnotesize] at (7,3.0) {Écran};
  % sources
  \coordinate (O)  at (0,0);
  \coordinate (S2) at (0,1.1);
  \coordinate (S1) at (0,-1.1);
  \fill[black] (S1) circle (2.4pt) node[left]{$S_1$};
  \fill[black] (S2) circle (2.4pt) node[left]{$S_2$};
  \node[left] at (O) {$O$};
  % point M sur l'écran
  \coordinate (M) at (7,2.3);
  \fill[black] (M) circle (2.4pt) node[above]{$M$};
  % axe central
  \draw[dashed,black!55] (O)--(7,0);
  % rayons d1, d2
  \draw[black,line width=0.9pt] (S1)--(M); \node at (3.2,0.95){$d_1$};
  \draw[black,line width=0.9pt] (S2)--(M); \node at (3.6,2.0){$d_2$};
  % cote a (entre les sources)
  \draw[<->,>=Stealth,black] (-0.55,-1.1)--(-0.55,1.1)
        node[midway,left]{$a$};
  % cote D (distance plan-écran)
  \draw[<->,>=Stealth,black] (0,-2.4)--(7,-2.4) node[midway,below]{$D$};
  % cote x (hauteur de M)
  \draw[<->,>=Stealth,black] (7.4,0)--(7.4,2.3) node[midway,right]{$x$};
\end{tikzpicture}
\caption{Dispositif des fentes de Young. La différence de marche
$\delta = d_1 - d_2 = a\,x/D$ fixe l'intensité au point $M$ situé à la hauteur $x$ sur l'écran.}
\label{fig:young}
\end{figure}

\begin{formule}[interfrange et différence de marche]
\[ \boxed{\,i = \frac{\lambda\,D}{a}\,}\qquad\qquad
   \boxed{\,d_1 - d_2 = \frac{a\,x}{D}\,} \]
\begin{ou}
  \item $i$ : \textbf{interfrange}, distance entre deux franges \emph{consécutives} de même
        nature, en \si{\metre} ;
  \item $\lambda$ : longueur d'onde, en \si{\metre} ;
  \item $a$ : distance entre les deux sources (ou les deux fentes), en \si{\metre} ;
  \item $D$ : distance entre le plan des sources et l'écran, en \si{\metre}, avec $D \gg a$ ;
  \item $x$ : position du point $M$ sur l'écran (mesurée depuis l'axe central), en \si{\metre}.
\end{ou}
\textbf{Lecture.} L'interfrange est \emph{proportionnel} à $\lambda$ et à $D$, et
\emph{inversement proportionnel} à $a$ : pour bien séparer les franges, on prend des fentes
très rapprochées ($a$ petit) et un écran éloigné ($D$ grand).
\end{formule}

\begin{exemple}[fentes de Young — calcul de l'interfrange]
Une lumière monochromatique de longueur d'onde $\lambda = \SI{600}{\nano\metre}$ (rouge)
éclaire deux fentes distantes de $a = \SI{0.5}{\milli\metre}$. L'écran est placé à
$D = \SI{2}{\metre}$. On veut l'interfrange $i$ et la position de la 3\textsuperscript{e} frange
brillante.

\smallskip
\textbf{Conversion en unités SI.} \quad
$\lambda = \SI{600}{\nano\metre} = \num{600e-9}\,\si{\metre} = \num{6e-7}\,\si{\metre}$,\quad
$a = \SI{0.5}{\milli\metre} = \num{5e-4}\,\si{\metre}$.

\smallskip
\textbf{Interfrange.}
\[ i = \frac{\lambda D}{a}
   = \frac{\num{6e-7}\times 2}{\num{5e-4}}
   = \frac{\num{1.2e-6}}{\num{5e-4}}
   = \num{2.4e-3}\,\si{\metre} = \SI{2.4}{\milli\metre}. \]
Deux franges brillantes voisines sont donc séparées de \SI{2.4}{\milli\metre} sur l'écran.

\smallskip
\textbf{Position de la 3\textsuperscript{e} frange brillante.} Les franges brillantes sont en
$x_n = n\,i$ (interférence constructive, $d_1 - d_2 = n\lambda$), donc :
\[ x_3 = 3\,i = 3\times \SI{2.4}{\milli\metre} = \SI{7.2}{\milli\metre}. \]

\smallskip
\textbf{Vérification par la différence de marche.} À la première frange brillante ($x = i$) :
\[ d_1 - d_2 = \frac{a\,x}{D} = \frac{\num{5e-4}\times \num{2.4e-3}}{2}
   = \num{6e-7}\,\si{\metre} = \lambda. \]
On retrouve bien $d_1 - d_2 = 1\times\lambda$ : la condition d'interférence \emph{constructive}
($n=1$) est satisfaite. Tout est cohérent.
\end{exemple}

% =====================================================================
\section{Synthèse : l'essentiel en un coup d'œil}
% =====================================================================

\begin{center}
\renewcommand{\arraystretch}{1.45}
\begin{tabular}{>{\raggedright\arraybackslash}m{4.7cm}
                >{\centering\arraybackslash}m{5.4cm}
                >{\centering\arraybackslash}m{3.3cm}}
\rowcolor{primary!18}
\textbf{Grandeur / relation} & \textbf{Formule} & \textbf{Unité} \\
\rowcolor{primary!5}
Période $\leftrightarrow$ fréquence & $T = \dfrac{1}{f}$ & \si{\second},\ \si{\hertz} \\
Longueur d'onde & $\lambda = vT = \dfrac{v}{f}$ & \si{\metre} \\
\rowcolor{primary!5}
Célérité & $v = \lambda f = \dfrac{\omega}{k}$ & \si{\metre\per\second} \\
Pulsation, nombre d'onde & $\omega = 2\pi f,\quad k = \dfrac{2\pi}{\lambda}$
   & \si{\radian\per\second},\ \si{\radian\per\metre} \\
\rowcolor{primary!5}
Élongation sinusoïdale & $y = A\sin(\omega t - kx + \varphi)$ & \si{\metre} \\
Vitesse / accélération & $v_{\mathrm P}=A\omega\cos(\cdot),\ a_{\mathrm P}=-\omega^2 y$
   & \si{\metre\per\second},\ \si{\metre\per\second\squared} \\
\rowcolor{primary!5}
Harmoniques (2 bouts fixes) & $\lambda_n = \dfrac{2L}{n},\ f_n=\dfrac{nc}{2L}$
   & \si{\metre},\ \si{\hertz} \\
Harmoniques (1 bout libre) & $\lambda_n = \dfrac{4L}{2n-1}$ & \si{\metre} \\
\rowcolor{primary!5}
Interférence constructive & $d_1-d_2 = n\lambda$ & \si{\metre} \\
Interférence destructive & $d_1-d_2 = (2n+1)\dfrac{\lambda}{2}$ & \si{\metre} \\
\rowcolor{primary!5}
Interfrange (Young) & $i = \dfrac{\lambda D}{a}$ & \si{\metre} \\
\end{tabular}
\end{center}

\vspace{0.3cm}
\begin{tcolorbox}[boxbase,colback=primarydark!6,colframe=primarydark,
   title={\textsc{À retenir absolument}}]
\begin{itemize}[leftmargin=1.5em,topsep=2pt,itemsep=2pt]
  \item Une onde transporte de l'\textbf{énergie}, jamais de la matière ; chaque point
        \emph{oscille sur place}.
  \item \textbf{Longitudinale} = vibration parallèle à la propagation (son) ;
        \textbf{transversale} = vibration perpendiculaire (corde, lumière).
  \item La \textbf{fréquence} est imposée par la source ; la \textbf{célérité} et la
        \textbf{longueur d'onde} dépendent du milieu.
  \item \textbf{Concordance de phase} : $d=n\lambda$ ; \textbf{opposition de phase} :
        $d=(2n+1)\lambda/2$. Ces mêmes conditions gouvernent les interférences.
  \item Une \textbf{onde stationnaire} (incidente + réfléchie) possède des \textbf{nœuds}
        fixes et des \textbf{ventres} ; une corde fixée ne vibre qu'à ses fréquences propres
        $f_n = nc/2L$.
\end{itemize}
\end{tcolorbox}

\end{document}
